\section{Coin Combinations I}
\label{sec:cses-coin-combinations-i}

\problemstatement{Given $n$ coins and a target $x$, count the number of
\emph{ordered} ways to make $x$ (a coin may be used many times, and
sequences differing only in order are counted separately). Output modulo
$10^9+7$.}

\begin{patternbox}
\emph{``Count ways, order matters''} with reusable coins is Dice
Combinations generalised to an arbitrary coin set. Let $\textit{dp}[s]$
count ordered sequences summing to $s$; the last coin is any $c$, so
$\textit{dp}[s]=\sum_c \textit{dp}[s-c]$.
\end{patternbox}

\subsection*{State, transition, and loop order}

$\textit{dp}[s]$ counts ordered sequences summing to $s$,
$\textit{dp}[0]=1$. The last coin used has some value $c$; the prefix
before it sums to $s-c$:
\[
  \textit{dp}[s]=\sum_{c:\,c\le s}\textit{dp}[s-c].
\]
The crucial detail is the \textbf{loop order}: the \emph{sum} $s$ runs on
the outside and coins on the inside. This way every coin is considered as
a possible \emph{last} move for each sum, so $(2\text{ then }3)$ and
$(3\text{ then }2)$ are both counted -- exactly what ``order matters''
requires.

\begin{ideabox}[Sum-Outer Loop Counts Permutations]
Putting the sum on the outer loop lets each sum choose \emph{any} coin as
its last step, generating all orderings. Reversing the loops (coins
outer) would count each multiset once -- the distinction between this
problem and Coin Combinations II (Section~\ref{sec:cses-coin-combinations-ii}).
\end{ideabox}

\subsection*{Algorithm}

\begin{enumerate}
  \item $\textit{dp}[0]=1$.
  \item For $s=1\ldots x$ (outer): for each coin $c\le s$ (inner),
        $\textit{dp}[s]\mathrel{+}=\textit{dp}[s-c]$, modulo $10^9+7$.
  \item Answer $\textit{dp}[x]$.
\end{enumerate}

\subsection*{Dry run}

\dryrun{Coins $\{2,3,5\}$, $x=9$; partial values.}

\begin{itemize}
  \item $\textit{dp}[0]=1$, $\textit{dp}[2]=1$, $\textit{dp}[3]=1$,
        $\textit{dp}[5]=\textit{dp}[3]+\textit{dp}[2]+\textit{dp}[0]=3$
        ($2{+}3$, $3{+}2$, $5$).
  \item Building up, $\textit{dp}[9]=8$ -- all ordered sequences of
        $\{2,3,5\}$ summing to $9$.
\end{itemize}

\subsection*{C++ implementation}

\begin{lstlisting}
#include <bits/stdc++.h>
using namespace std;
const int MOD = 1e9 + 7;

int main() {
    int n, x;
    cin >> n >> x;
    vector<int> coins(n);
    for (int& c : coins) cin >> c;

    vector<long long> dp(x + 1, 0);
    dp[0] = 1;
    for (int s = 1; s <= x; s++)          // sum outer -> permutations
        for (int c : coins)
            if (c <= s) dp[s] = (dp[s] + dp[s - c]) % MOD;

    cout << dp[x] << "\n";
    return 0;
}
\end{lstlisting}

\begin{complexitybox}
\textbf{Time Complexity.} $O(nx)$. \textbf{Space Complexity.} $O(x)$.
\end{complexitybox}

\begin{takeawaybox}
Ordered coin counting = sum-outer, coin-inner loops. Memorise this loop
order as the switch that toggles between counting permutations (here) and
combinations (next section). Getting it backward is the single most common
DP counting bug.
\end{takeawaybox}
